% GENERATED FILE. Edit canonical lessons, then run npm run generate:books.
% canonical-source-hash: fnv1a64:2f28e776e4ad4a64
% canonical-lessons: FR-C07-lune, FR-C07-lundi, FR-C07-mardi, FR-C07-mercredi, FR-C07-jeudi, FR-C07-vendredi, FR-C07-samedi, FR-C07-dimanche, FR-C07-practice

\chapter{The Days of the Week}
\label{ch:days}
\hlchaptermodality{7}

\begin{chapteropening}
\textbf{By the end of this chapter:} \emph{I can name the days of the week in French and use them in simple plans.}

Seven day-names that are a map of the sky. Five carry a Roman planet-god and all of them carry Latin dies, \textquotedblleft{}day\textquotedblright{} -- so half of every weekday comes free, and English names the same gods on the same days under Germanic names. Then two days where a later faith wrote over the planets entirely.
\end{chapteropening}

\section[la lune]{la lune — the Moon}

\label{lesson:FR-C07-lune}



\begin{quote}
The seven-day week is one of the oldest things any language carries, and the French day-names are a \textbf{map of the sky}. Before we read the map, one of the objects on it.
\end{quote}

\begin{sounds}
\begin{itemize}
\raggedright
  \item \texttt{front-rounded-u} — the French \textbf{u} has no English equivalent. Say \textquotedblleft{}ee,\textquotedblright{} and then, without moving your tongue, round your lips as though for \textquotedblleft{}oo.\textquotedblright{} That is \textbf{u}. So \textbf{lune} is \emph{lün}, never \emph{loon}.
  \item \texttt{silent-final-e} — the closing \textbf{-e} is a whisper, not a syllable of its own.
\end{itemize}
\end{sounds}

\begin{cousinweb}
\textbf{la lune} is feminine, so it takes \textbf{la}, the article you already know.

It comes from Latin \textbf{lūna}, which English kept in a small family:

\begin{itemize}
\raggedright
  \item \textbf{lunar} — of the Moon
  \item \textbf{lunation} — one full cycle of it
  \item \textbf{lunatic} — someone the Moon was believed to have unsettled, back when moonlight was thought to disturb the mind
\end{itemize}
\end{cousinweb}

\subsection*{Your turn}
Say these aloud:

\begin{itemize}
\raggedright
  \item \textquotedblleft{}lune\textquotedblright{} — \emph{lün}, lips rounded and tongue forward
  \item \textquotedblleft{}la lune\textquotedblright{} — feminine, so \emph{la}
  \item \textquotedblleft{}lune\textquotedblright{} then English \textquotedblleft{}lunar, lunatic\textquotedblright{} — the same \emph{lūna}
\end{itemize}

\begin{tcolorbox}[breakable,colback=teal!4,colframe=teal!35!black,title={Before you move on}]
How do you say the Moon? (\textbf{la lune}, \emph{la lün}.) Which article does it take, and why? (\textbf{la} — the word is feminine.) How do you make the French \textbf{u}? (Say \textquotedblleft{}ee,\textquotedblright{} then round the lips without moving the tongue.) Next: the day named after it.
\end{tcolorbox}
\section[lundi]{lundi — Monday}

\label{lesson:FR-C07-lundi}



\begin{quote}
You have the Moon. Add one ending and you have the first day of the week — and, with it, the second half of five more words.
\end{quote}

\begin{sounds}
\begin{itemize}
\raggedright
  \item \textbf{lundi} = \emph{luh(n)-DEE}, opening with the nasal \emph{un} you met in the number \textbf{un}.
  \item The weight falls on the \textbf{second} syllable. French leans on the end of a word where English leans near the front: \emph{MON-day} against \emph{lun-DI}.
\end{itemize}
\end{sounds}

\begin{grammarlens}[title={Grammar Lens: every weekday ends in \textquotedblleft{}day\textquotedblright{}}]
Every French weekday ends in \textbf{-di}, always said \emph{DEE}. That \textbf{-di} is Latin \textbf{diēs}, \textquotedblleft{}day\textquotedblright{} — the same word English puts at the \emph{end} of \emph{Monday}, and French puts at the end too.

So a French weekday reads as \textbf{\textquotedblleft{}{[}god{]}'s day,\textquotedblright{}} the god in front and the day behind:

\begin{quote}
\textbf{lundi} = \emph{lūnae} (\textquotedblleft{}of the Moon\textquotedblright{}) + \emph{diēs} (\textquotedblleft{}day\textquotedblright{}) = \textbf{the Moon's day}
\end{quote}

The Romans named each day for a planet-god, and the planets were the gods. Learn to hear that \textbf{-di} and half of five more words is already yours.
\end{grammarlens}

\begin{culture}
English \emph{Monday} says the identical thing — \emph{Moon} plus \emph{day} — but out of Germanic parts rather than Latin ones. The two languages did not copy each other. They inherited the \textbf{same Roman week} and each translated it into its own words. That is why the days will keep lining up as we go, and why they will line up through their \emph{meanings} rather than their sounds.
\end{culture}

\subsection*{Your turn}
Say these aloud:

\begin{itemize}
\raggedright
  \item \textquotedblleft{}lundi\textquotedblright{} — \emph{luh(n)-DEE}, weight on the end
  \item it split — \textquotedblleft{}lune\textquotedblright{} plus \textquotedblleft{}-di\textquotedblright{}: the Moon's day
  \item \textquotedblleft{}lundi, Monday\textquotedblright{} — the same sentence in two languages
\end{itemize}

\begin{tcolorbox}[breakable,colback=teal!4,colframe=teal!35!black,title={Before you move on}]
How do you say Monday? (\textbf{lundi}, \emph{luh(n)-DEE}.) What does the \textbf{-di} mean, and from which Latin word? (\textquotedblleft{}Day,\textquotedblright{} from \textbf{diēs}.) What does \emph{lundi} mean once you take it apart, and what does \emph{Monday} mean? (Both: \textbf{the Moon's day}.) Next: the day of the war-god.
\end{tcolorbox}
\section[mardi]{mardi — Tuesday}

\label{lesson:FR-C07-mardi}



\begin{quote}
The second day carries the best story in the week: \emph{mardi} and \emph{Tuesday} look nothing alike and are, underneath, exactly the same word.
\end{quote}

\begin{sounds}
\begin{itemize}
\raggedright
  \item \textbf{mardi} = \emph{mar-DEE}, with the \textbf{r} made at the back of the throat, as in \textbf{quatre}.
  \item The familiar \textbf{-di} closes it, so only the first syllable is new.
\end{itemize}
\end{sounds}

\begin{cousinweb}
\textbf{mardi} is \textbf{Martis diēs}, \textquotedblleft{}the day of \textbf{Mars}\textquotedblright{} — the god of war, and the red planet named for him.
\end{cousinweb}

\begin{culture}
English \emph{Tuesday} is named for \textbf{Tiw}, the Germanic war-god. French \emph{mardi} is named for \textbf{Mars}, the Roman one. Different gods, different sounds — and the \textbf{same day}, because it is the same slot in the same week.

Here is how that happened. The Romans mapped their gods onto the planets and named the days for them. When the week reached the Germanic-speaking peoples, they did not borrow the Roman names; they \textbf{swapped in their own gods}, matching them role for role. Scholars call that \emph{interpretatio germanica}: translating a pantheon rather than a word.

So the two languages will keep agreeing about \textbf{which} god each day belongs to while disagreeing about that god's name. Once you know that, every remaining weekday has an English key.
\end{culture}

\subsection*{Your turn}
Say these aloud:

\begin{itemize}
\raggedright
  \item \textquotedblleft{}mardi\textquotedblright{} — \emph{mar-DEE}, the \emph{r} at the back of the throat
  \item \textquotedblleft{}lundi, mardi\textquotedblright{} — the first two, both ending \emph{DEE}
  \item the bridge — \textquotedblleft{}mardi is Mars, Tuesday is Tiw, both the war-god\textquotedblright{}
\end{itemize}

\begin{tcolorbox}[breakable,colback=teal!4,colframe=teal!35!black,title={Before you move on}]
How do you say Tuesday? (\textbf{mardi}, \emph{mar-DEE}.) Which god is it named for? (\textbf{Mars}, the war-god.) Why do \emph{mardi} and \emph{Tuesday} line up despite sharing no sound? (Same day, same role — Latin used \emph{Mars}, the Germanic peoples put in their own war-god \textbf{Tiw}.) Next: the messenger's day.
\end{tcolorbox}
\section[mercredi]{mercredi — Wednesday}

\label{lesson:FR-C07-mercredi}



\begin{quote}
The longest day-name of the week, and the first place to try the key you were handed: find the god, and English will tell you the day.
\end{quote}

\begin{sounds}
\begin{itemize}
\raggedright
  \item \textbf{mercredi} = \emph{mair-kruh-DEE}, three syllables with the weight at the end.
  \item Two throat-\textbf{r} sounds in one word, one in each of the first two syllables. Take it slowly; this word is worth drilling on its own.
\end{itemize}
\end{sounds}

\begin{cousinweb}
\textbf{mercredi} is \textbf{Mercuriī diēs}, \textquotedblleft{}the day of \textbf{Mercury}\textquotedblright{} — the messenger of the gods, and the swiftest planet. English kept the god's name in \textbf{mercurial}, \textquotedblleft{}quick-changing,\textquotedblright{} and gave it to the metal that runs like water.
\end{cousinweb}

\begin{culture}
Apply the key from \emph{mardi}. English \emph{Wednesday} is \textbf{Woden's day} — and Woden is the god the Romans matched to Mercury, both of them travellers, guides and carriers of the dead. Same day, same role, two names. English still spells the \emph{d} it stopped pronouncing.
\end{culture}

\subsection*{Your turn}
Say these aloud:

\begin{itemize}
\raggedright
  \item \textquotedblleft{}mercredi\textquotedblright{} — \emph{mair-kruh-DEE}, slowly, both \emph{r} sounds at the back
  \item \textquotedblleft{}lundi, mardi, mercredi\textquotedblright{} — three down
  \item the key again — \textquotedblleft{}Mercury for the Romans, Woden for the Germanic peoples, one day\textquotedblright{}
\end{itemize}

\begin{tcolorbox}[breakable,colback=teal!4,colframe=teal!35!black,title={Before you move on}]
How do you say Wednesday? (\textbf{mercredi}, \emph{mair-kruh-DEE}.) Which god is it named for? (\textbf{Mercury}, the messenger.) Which god does English name the same day for? (\textbf{Woden} — the god Rome matched to Mercury.) Next: the day of the king of the gods.
\end{tcolorbox}
\section[jeudi]{jeudi — Thursday}

\label{lesson:FR-C07-jeudi}



\begin{quote}
The fourth day belongs to the king of the gods, and it introduces the sound most people get wrong on their first attempt at French.
\end{quote}

\begin{sounds}
\begin{itemize}
\raggedright
  \item \texttt{j-zh} — the French \textbf{j} is the soft buzz in the middle of English \emph{mea\textbf{s}ure} or \emph{vi\textbf{si}on}. It is never the hard \emph{j} of \emph{jam}.
  \item \texttt{vowel-eu} — the rounded \emph{eu} you learned in \textbf{deux}.
  \item Together: \textbf{jeudi} = \emph{zheu-DEE}.
\end{itemize}
\end{sounds}

\begin{cousinweb}
\textbf{jeudi} is \textbf{Jovis diēs}, \textquotedblleft{}the day of \textbf{Jove}\textquotedblright{} — Jove being the other name of \textbf{Jupiter}, king of the gods and the largest planet. English keeps \emph{Jove} in the phrase \emph{by Jove}, and keeps the god's temperament in \textbf{jovial}: cheerful, the way people born under Jupiter were said to be.
\end{cousinweb}

\begin{culture}
The key once more. English \emph{Thursday} is \textbf{Thor's day} — Thor, the thunder-god, matched to Jupiter, who also threw the lightning. Two pantheons, one sky-king, one day. Every weekday so far has agreed on the god's \emph{job} and disagreed on the god's \emph{name}.
\end{culture}

\subsection*{Your turn}
Say these aloud:

\begin{itemize}
\raggedright
  \item \textquotedblleft{}jeudi\textquotedblright{} — \emph{zheu-DEE}, the soft \emph{j} of \textquotedblleft{}measure\textquotedblright{}
  \item \textquotedblleft{}mercredi, jeudi\textquotedblright{} — the pair
  \item the bridge — \textquotedblleft{}Jupiter and Thor, both the thunder-king, one day\textquotedblright{}
\end{itemize}

\begin{tcolorbox}[breakable,colback=teal!4,colframe=teal!35!black,title={Before you move on}]
How do you say Thursday? (\textbf{jeudi}, \emph{zheu-DEE}.) What English sound is the French \textbf{j}? (The buzz in \emph{measure}, never the \emph{j} of \emph{jam}.) Which god is \emph{jeudi} named for, and which god does \emph{Thursday} name? (\textbf{Jupiter}, or Jove — and \textbf{Thor}.) Next: the day of the morning star.
\end{tcolorbox}
\section[vendredi]{vendredi — Friday}

\label{lesson:FR-C07-vendredi}



\begin{quote}
The fifth day closes the planet-week, and it is the only one named for a goddess.
\end{quote}

\begin{sounds}
\begin{itemize}
\raggedright
  \item \texttt{nasal-en} — the opening \textbf{en} is nasal, sent through the nose like the \emph{un} of \textbf{lundi} but with the mouth more open. So \textbf{vendredi} = \emph{van-druh-DEE}.
  \item The \textbf{r} is the throat-\textbf{r} again, and the ending is the familiar \emph{DEE}.
\end{itemize}
\end{sounds}

\begin{cousinweb}
\textbf{vendredi} is \textbf{Veneris diēs}, \textquotedblleft{}the day of \textbf{Venus}\textquotedblright{} — goddess of love and the brightest thing in the sky after the Sun and Moon. Her name is still in English \textbf{venerate}, \textquotedblleft{}to hold in reverence.\textquotedblright{}
\end{cousinweb}

\begin{culture}
The last application of the key, and the cleanest. English \emph{Friday} is \textbf{Frigg's day} — Frigg, the Germanic goddess of love and marriage, matched to Venus. Five days, five gods, and every single pairing made by role.

That is the whole planet-week: \textbf{lundi} the Moon, \textbf{mardi} Mars, \textbf{mercredi} Mercury, \textbf{jeudi} Jupiter, \textbf{vendredi} Venus. Two days are left, and neither of them is a planet — which is the next surprise.
\end{culture}

\subsection*{Your turn}
Say these aloud:

\begin{itemize}
\raggedright
  \item \textquotedblleft{}vendredi\textquotedblright{} — \emph{van-druh-DEE}, the \emph{en} through the nose
  \item the work week — \textquotedblleft{}lundi, mardi, mercredi, jeudi, vendredi\textquotedblright{}
  \item each as a god — \textquotedblleft{}Moon, Mars, Mercury, Jupiter, Venus\textquotedblright{}
\end{itemize}

\emph{Twice through:} the five weekdays, in order

\begin{tcolorbox}[breakable,colback=teal!4,colframe=teal!35!black,title={Before you move on}]
How do you say Friday? (\textbf{vendredi}, \emph{van-druh-DEE}.) Which goddess is it named for, and who does English name it for? (\textbf{Venus} — and \textbf{Frigg}.) Give the five planet-days in order. (\emph{Lundi, mardi, mercredi, jeudi, vendredi}.) Next: the two days where the planets stop.
\end{tcolorbox}
\section[samedi]{samedi — Saturday}

\label{lesson:FR-C07-samedi}



\begin{quote}
Five days named for planets — and then two that are not. The weekend is where \textbf{religion overwrote astronomy}, and French shows both layers.
\end{quote}

\begin{sounds}
\begin{itemize}
\raggedright
  \item \textbf{samedi} = \emph{sam-DEE}. The \textbf{a} is open, and the middle \textbf{e} is so light it nearly disappears: two syllables in practice, not three.
  \item \texttt{nasal-none} — the \textbf{am} here is \textbf{not} nasal, because the \emph{m} is followed by a vowel that pulls it into the next syllable.
\end{itemize}
\end{sounds}

\begin{cousinweb}
\textbf{samedi} still carries the \textbf{-di} of \emph{diēs}, so the shape of the week holds. What changed is the first half. It is \textbf{sambatī diēs}, from Latin \textbf{Sabbatum}, which is Hebrew \textbf{shabbāt} — \textquotedblleft{}to cease, to rest.\textquotedblright{}

Not a god. An instruction.
\end{cousinweb}

\begin{culture}
By the Roman pattern this should have been \textbf{Saturn's day} — and in English it still is. \emph{Saturday} is the one weekday English kept a Roman god for, while replacing the other four with Germanic ones.

French did the reverse. It kept the Roman gods on the weekdays and let this one go: Latin-speaking Christians took the Jewish \textbf{Sabbath} and put it here, so the planet was written over by the rest-day.

So \emph{Saturday} and \emph{samedi} are the same day under two ideas — \textbf{the planet against the Sabbath} — and the two languages swapped exactly which half of the week they preserved.
\end{culture}

\subsection*{Your turn}
Say these aloud:

\begin{itemize}
\raggedright
  \item \textquotedblleft{}samedi\textquotedblright{} — \emph{sam-DEE}, the middle \emph{e} almost gone
  \item \textquotedblleft{}vendredi, samedi\textquotedblright{} — the pair
  \item the contrast — \textquotedblleft{}samedi is the Sabbath, Saturday is Saturn\textquotedblright{}
\end{itemize}

\begin{tcolorbox}[breakable,colback=teal!4,colframe=teal!35!black,title={Before you move on}]
How do you say Saturday? (\textbf{samedi}, \emph{sam-DEE}.) Where does the first half come from? (The \textbf{Sabbath} — Latin \emph{Sabbatum}, Hebrew \emph{shabbāt}, \textquotedblleft{}to rest.\textquotedblright{}) Which planet-god does English keep on this day that French dropped? (\textbf{Saturn}.) Next: the last day, where even the \textbf{-di} moves.
\end{tcolorbox}
\section[dimanche]{dimanche — Sunday}

\label{lesson:FR-C07-dimanche}



\begin{quote}
Six days have ended in \textbf{-di}. The seventh breaks the pattern in public — and once you see how, it is not a break at all.
\end{quote}

\begin{sounds}
\begin{itemize}
\raggedright
  \item \texttt{ch-sh} — the French \textbf{ch} is \emph{sh}, as in English \emph{machine}. Never \emph{ch} as in \emph{church}.
  \item \texttt{nasal-an} — the \textbf{an} is nasal, the same open nasal as the \textbf{en} of \textbf{vendredi}: French spells that one sound two ways.
  \item Together: \textbf{dimanche} = \emph{dee-MAHNSH}.
\end{itemize}
\end{sounds}

\begin{grammarlens}[title={Grammar Lens: the ending that moved to the front}]
Look at the first two letters. \textbf{Di}manche. That is the \textbf{same \emph{diēs}} every other day of the week ends with, sitting at the front instead.

\begin{quote}
\textbf{diēs Dominica} $\to$ \emph{di-manche} $\to$ \textbf{dimanche}
\end{quote}

Latin let the two words come in either order, and for this one day the order that survived is \textquotedblleft{}day of the Lord\textquotedblright{} rather than \textquotedblleft{}the Lord's day.\textquotedblright{} So the pattern did not break; it flipped. All seven days still carry \emph{diēs}, and \emph{dimanche} is the only one that wears it where you can see it first.
\end{grammarlens}

\begin{culture}
\emph{Dominica} comes from \textbf{Dominus}, \textquotedblleft{}lord, master\textquotedblright{} — which English holds in \textbf{dominion}, \textbf{dominate}, and, worn much smoother, \textbf{dame}.

English kept \textbf{Sun}day, the star. French dropped it entirely for \textquotedblleft{}the Lord's day.\textquotedblright{} Put beside \emph{samedi}, that makes the weekend two different rewrites of the same Roman week:

\begin{itemize}
\raggedright
  \item \textbf{samedi} — the planet replaced by the \textbf{Sabbath}, a day of rest
  \item \textbf{dimanche} — the Sun replaced by the \textbf{Lord}, a day of worship
\end{itemize}

Five days of Roman sky, then two days where a later faith wrote over it. The French week is a record of both.
\end{culture}

\subsection*{Your turn}
Say these aloud:

\begin{itemize}
\raggedright
  \item \textquotedblleft{}dimanche\textquotedblright{} — \emph{dee-MAHNSH}, the \emph{ch} as \emph{sh}
  \item \textquotedblleft{}samedi, dimanche\textquotedblright{} — the weekend
  \item the whole week — \textquotedblleft{}lundi, mardi, mercredi, jeudi, vendredi, samedi, dimanche\textquotedblright{}
\end{itemize}

\begin{tcolorbox}[breakable,colback=teal!4,colframe=teal!35!black,title={Before you move on}]
How do you say Sunday? (\textbf{dimanche}, \emph{dee-MAHNSH}.) Where did its \emph{diēs} go? (\textbf{To the front} — \emph{diēs Dominica}.) What does it mean literally? (\textquotedblleft{}The \textbf{Lord's} day.\textquotedblright{}) What does English keep on this day that French dropped? (The \textbf{Sun}.) Next: the whole week, run together.
\end{tcolorbox}
\section[Practice]{Practice — the whole week}

\label{lesson:FR-C07-practice}



\begin{quote}
Seven days, taken one at a time. Now put them in a row and make the row automatic — a week you have to assemble is a week you cannot use.
\end{quote}

\subsection*{Your turn: the seven, in order}
Read down the column, then cover the right-hand side and read down the left:

\noindent
\begin{tabularx}{\linewidth}{@{}>{\raggedright\arraybackslash}X>{\raggedright\arraybackslash}X@{}}
\toprule
\textbf{French} & \textbf{English} \\
\midrule
\textbf{lundi} & Monday \\
\textbf{mardi} & Tuesday \\
\textbf{mercredi} & Wednesday \\
\textbf{jeudi} & Thursday \\
\textbf{vendredi} & Friday \\
\textbf{samedi} & Saturday \\
\textbf{dimanche} & Sunday \\
\bottomrule
\end{tabularx}

\emph{Twice through:} the seven days, in order

Then start from \textbf{jeudi} and go round to \emph{mercredi}. A list you can only say from the top is a list you have memorised rather than learned.

\subsection*{Your turn: name the god}
Five of the seven still name a planet. Say each day, then what is inside it:

\begin{itemize}
\raggedright
  \item \textbf{lundi} — \emph{la lune}, the Moon
  \item \textbf{mardi} — Mars
  \item \textbf{mercredi} — Mercury
  \item \textbf{jeudi} — Jupiter
  \item \textbf{vendredi} — Venus
\end{itemize}

Every one of those ends in \textbf{-di}, Latin \emph{diēs}, \textquotedblleft{}day.\textquotedblright{}

\begin{grammarlens}[title={What you've built this chapter}]
\begin{itemize}
\raggedright
  \item \textbf{the -di rule} — every weekday ends in Latin \emph{diēs}, \textquotedblleft{}day,\textquotedblright{} so half of five words came free.
  \item \textbf{the planet week} — Moon, Mars, Mercury, Jupiter, Venus, named by Rome.
  \item \textbf{the key to English} — \emph{Tuesday}, \emph{Wednesday}, \emph{Thursday} and \emph{Friday} name the \textbf{same} gods by their Germanic names, matched role for role. The days agree about the god's job and disagree about the god's name.
  \item \textbf{the weekend rewrite} — \emph{samedi} is the \textbf{Sabbath} where English keeps Saturn; \emph{dimanche} is the \textbf{Lord's day} where English keeps the Sun. And in \emph{dimanche} the \emph{diēs} moved to the front.
\end{itemize}
\end{grammarlens}

\begin{tcolorbox}[breakable,colback=teal!4,colframe=teal!35!black,title={Before you move on}]
Give the seven days without stopping. Now start from \textbf{jeudi}. What does the \textbf{-di} mean? (\textquotedblleft{}Day,\textquotedblright{} Latin \emph{diēs}.) Which two days are not named for a planet, and what replaced them? (\textbf{samedi}, the Sabbath; \textbf{dimanche}, the Lord's day.) Why do \emph{mardi} and \emph{Tuesday} mean the same thing? (Same day, same role — the war-god, called \textbf{Mars} in Latin and \textbf{Tiw} in Germanic.)
\end{tcolorbox}
